% === Krasnokutskaya positioning block — VERIFIED from PDF (v2 post-review) ===
This paper is closest in spirit to the structural auction literature on
small-business preferences in procurement, notably
\citet{krasnokutskayaseim2011} and \citet{krasnokutskaya2011}. Three
differences sharpen the contribution:

\textit{Policy instrument.}
\citet{krasnokutskayaseim2011} study the California Small Business
Preference program: a \textbf{5\% bid discount} for small firms in
Caltrans highway auctions; small firms win if their bid is within 5\%
of the lowest nonfavored bid. The sample is 697 projects, 3{,}034 bids,
2002--2005, with an average of 4.37 bidders per project (1.75 small
+ 2.62 large). Our setting is a \textbf{complete restriction}
(SME-only tendering) in Brazilian medical-supply procurement. Complete
restriction is an orders-of-magnitude stronger intervention: the
non-SME extensive margin is entirely eliminated rather than tilted.

\textit{Primitive-invariance identification.}
The March 2018 regime break in our setting excludes non-SMEs from
Group-65 auctions ex post. Non-SME cost primitives are unaffected;
we test and confirm primitive invariance directly
(KS $D \leq 0.04$, mean shift $< 0.02$ of reference price).
\citet{krasnokutskayaseim2011} do not have a comparable regime break;
type-primitive invariance is an assumption, not a test.

\textit{Welfare magnitudes.}
\citet{krasnokutskayaseim2011}'s policy-relevant estimate (Table 11)
finds that the California 5\% Small Business Preference raises
procurement cost by \textbf{0.5\%} relative to no discount and shifts
small-firm share of work from 12.5\% to 15.6\%. Counterfactual
discounts of 25\% and 50\% raise cost by 4.1\% and 7.8\% respectively.
Our SME-only set-aside generates a 10.4\% fiscal cost, 11.8\% DWL of
procurement value (R\$295k over 18 months for Convite G65 alone),
roughly \textbf{20$\times$} the baseline K\&S 5\% estimate and
comparable in magnitude to their 50\% discount counterfactual. The
proportionality is consistent with moving from a mild preference to
a full restriction.

\textit{Cost asymmetries.}
\citet{krasnokutskayaseim2011} estimate that small firms bid 5.4\%
(parametric) to 8.1\% (OLS) higher than large firms on comparable
projects. Our Preg\~ao Haile-Tamer upper bound puts the SME-vs-non-SME
bid gap at about 17\% (the gap is an upper bound on the primitive cost
gap). Under Convite CPV identification, the gap at clean type
identification is near zero; under a stricter canonical classification
it widens to 5.4\%.

\textit{Unobserved heterogeneity and cost decomposition.}
\citet{krasnokutskaya2011} shows that 34.4\% of the variation in
bidders' costs in Michigan highway procurement is attributable to
private information, the remaining 66\% to common or unobserved
auction-level heterogeneity. Our reference-price normalization absorbs
part of the auction-level heterogeneity; residual variation is more
directly interpretable as cost-primitive variation across firms.
Estimating without the ref-price normalization, and comparing the two
deconvolution patterns, is natural future work.
